\documentclass[11pt]{article}

\usepackage[margin=1in]{geometry}
\usepackage{amsmath, amssymb, amsthm}
\usepackage{bm}
\usepackage{hyperref}
\usepackage{enumitem}
\usepackage{booktabs}

\newcommand{\R}{\mathbb{R}}
\newcommand{\E}{\mathbb{E}}
\newcommand{\Prob}{\mathbb{P}}
\newcommand{\Normal}{\mathcal{N}}
\newcommand{\dt}{\Delta t}
\newcommand{\bs}{\mathbf s}
\newcommand{\bz}{\mathbf z}
\newcommand{\by}{\mathbf y}

\title{Mathematical Formulation of Learning and Sampling for\\ Extreme Ground-Motion Events}
\author{ExtremeEventFlows}
\date{}

\begin{document}
\maketitle

\tableofcontents

\section{Overview}

This document gives the mathematical description of the physical simulation model,
the rare-event target density, the trainable autoregressive normalizing-flow proposal,
and the estimators used to learn that proposal and to estimate rare-event
probabilities. It mirrors the implementation in \texttt{extreme\_event\_flows/}:
\texttt{source.py}, \texttt{ground\_motion.py}, \texttt{application.py},
\texttt{target.py}, \texttt{flow.py}, \texttt{training.py}, and
\texttt{estimators.py}.

The overall goal is to estimate the probability of a rare, extreme event
(e.g. exceeding an intensity-measure threshold at a site over an exposure window)
under a physical seismic-source and ground-motion model, using importance
sampling with a learned proposal that concentrates mass on the failure region.

\section{The physical event sequence}

\subsection{Variable-length marked point process}

The seismic source is modeled as a variable-length sequence of events
$\mathcal X_N = (\bs_1,\dots,\bs_N)$ arriving on $[0,T]$, where $T$ is the
exposure time. Each event vector $\bs_i$ collects an arbitrary number of
physical source parameters, always including an inter-event time increment
$\dt_i$ so that the $i$-th event occurs at time
\[
t_i = \sum_{j=1}^i \dt_j .
\]
In the reference example,
\[
\bs_i = (\dt_i, M_i, x_i, y_i, d_i),
\]
i.e. time increment, magnitude, and 3-D location, but the framework supports
any named set of scalar event variables (\texttt{source.py:
GaussianLatentMarkovKernel}).

The process is Markov in the event history: the next event depends only on
the immediately preceding event. The joint density of a sequence with exactly
$N$ events inside $[0,T]$ is
\begin{equation}
p(\mathcal X_N)
=
\prod_{i=1}^{N} p(\bs_i \mid \bs_{i-1})\;
\Prob\big(\dt_{N+1} > T - t_N \mid \bs_N\big),
\label{eq:joint-source}
\end{equation}
where the last factor is the survival probability that no further event
occurs before the exposure horizon $T$ is reached, and $\bs_0$ denotes a
reference (initial) state.

\subsection{Latent-Gaussian transition kernel}

Each physical coordinate $s_{i,j}$ of event $i$ is related to a latent
Gaussian coordinate $z_{i,j}$ by an invertible scalar transform $T_j$,
\begin{equation}
s_{i,j} = T_j(z_{i,j}), \qquad j = 1,\dots,d,
\end{equation}
implemented as monotone bijections in \texttt{transforms.py} (e.g. identity,
softplus/exponential maps enforcing positivity or bounds, each with a
tractable Jacobian $\log \lvert T_j'(z)\rvert$).

Conditional on the previous physical event $\bs_{i-1}$, the latent vector
$\bz_i = (z_{i,1},\dots,z_{i,d})$ is jointly Gaussian,
\begin{equation}
\bz_i \mid \bs_{i-1} \sim \Normal(\bm\mu_i,\, LL^\top),
\qquad
\bm\mu_i = \bm\mu_0 + A\left(\frac{\bs_{i-1} - \bs_{\mathrm{ref}}}{\bm\sigma_{\mathrm{state}}}\right),
\label{eq:latent-transition}
\end{equation}
where $\bm\mu_0 \in \R^d$ is a base mean, $A \in \R^{d\times d}$ is a
transition matrix creating Markov dependence on every previous-event
coordinate, $L$ is a lower-triangular Cholesky factor giving a fully coupled
covariance $LL^\top$ across the coordinates of the current event, and
$\bs_{\mathrm{ref}}, \bm\sigma_{\mathrm{state}}$ are reference values and
scales that non-dimensionalize the conditioning feature.

The physical event density follows by the change-of-variables formula,
\begin{equation}
p(\bs_i \mid \bs_{i-1})
=
\Normal\!\big(\bz_i;\, \bm\mu_i, LL^\top\big)
\prod_{j=1}^d \left\lvert \frac{d z_{i,j}}{d s_{i,j}} \right\rvert ,
\qquad \bz_i = (T_1^{-1}(s_{i,1}),\dots,T_d^{-1}(s_{i,d})).
\label{eq:event-density}
\end{equation}

\subsection{Time-increment survival}

Because the time increment $\dt$ is one coordinate of $\bs_i$ with its own
marginal Gaussian latent coordinate (mean $\mu_i^{(\dt)}$, variance
$\sigma^2_{(\dt)} = \sum_k L_{(\dt),k}^2$ from row $(\dt)$ of $L$), the
survival probability that no event occurs in the remaining exposure window
$T - t_N$ is
\begin{equation}
\Prob(\dt_{N+1} > r \mid \bs_N)
=
\Phi\!\left(-\frac{T_{\dt}^{-1}(r) - \mu_{N+1}^{(\dt)}}{\sigma_{(\dt)}}\right),
\qquad r = T - t_N,
\label{eq:survival}
\end{equation}
with $\Phi$ the standard normal CDF (\texttt{log\_ndtr} in the code), and the
degenerate case $r \le$ the transform's lower bound returns log-survival
$0$.

Sampling the sequence (\texttt{source.py: sample\_sequence}) proceeds by
repeatedly drawing $\bz_i \sim \Normal(\bm\mu_i, LL^\top)$, mapping to
physical units, and accepting the event only while $t_i \le T$; the process
terminates the first time a proposed event time exceeds $T$.

\section{Ground-motion model}

\subsection{Feature construction}

For each accepted source event $i$ and a fixed site, a feature map
$\bm f_i = f(\bs_i, \text{site})$ is built (\texttt{ground\_motion.py:
GroundMotionFeatureBuilder}). The reference implementation uses
\begin{equation}
\bm f_i = \big(1,\; M_i,\; \log(R_i + R_0),\; \log(V_{s30}/V_{s30}^{\mathrm{ref}}),\; \dots \big),
\end{equation}
where $R_i = \lVert (x_i,y_i,d_i) - (x_{\mathrm{site}}, y_{\mathrm{site}}, d_{\mathrm{site}}) \rVert_2$
is the 3-D source-to-site distance, $R_0$ is a small offset avoiding a
singularity at zero distance, and additional source/site scalar fields may be
appended.

\subsection{Multi-output lognormal ground-motion model}

For $K$ intensity-measure outputs (e.g.\ PGA, $SA(0.2\,\mathrm{s})$,
$SA(1.0\,\mathrm{s})$), a linear model in log space with correlated residuals
gives, for event $i$,
\begin{equation}
\bm\mu_i^{(Y)} = B\bm f_i, \qquad
\ln \bm Y_i = \bm\mu_i^{(Y)} + L_Y \bm\epsilon_i, \qquad
\bm\epsilon_i \sim \Normal(0, I_K),
\label{eq:gmm}
\end{equation}
where $B \in \R^{K\times \dim(\bm f)}$ are per-output regression
coefficients and $L_Y$ is a lower-triangular Cholesky factor of the
within-event residual covariance, coupling the $K$ outputs (e.g. PGA and
spectral accelerations are correlated within the same event). The residual
vector $\bm\epsilon_i$ is a standard normal latent variable exactly as in a
normalizing flow, with density
\begin{equation}
p(\bm\epsilon_i) = \prod_{k=1}^K \Normal(\epsilon_{i,k}; 0, 1).
\label{eq:gmm-density}
\end{equation}
One latent residual is added per ground-motion output; the flow automatically
grows to include it (\texttt{ground\_motion.py:
MultiOutputGaussianGroundMotionModel}).

\section{Joint physical (``original'') density}

Composing the source and ground-motion models gives the joint density of a
complete simulated sample $x = (\mathcal X_N, \{\bm Y_i\}_{i=1}^N)$,
\begin{equation}
p(x) = \underbrace{p(\mathcal X_N)}_{\text{Eq.~\eqref{eq:joint-source}}}
\times
\underbrace{\prod_{i=1}^N p(\bm\epsilon_i)}_{\text{Eq.~\eqref{eq:gmm-density}}},
\label{eq:original-density}
\end{equation}
computed in \texttt{application.py: log\_original\_density} as the sum of a
source log-density and a ground-motion (residual) log-density.

\section{Rare-event target}

\subsection{Performance function}

A scalar performance function $g(x)$ maps a complete sample to a real number
such that the rare/failure event is $F = \{x : g(x) \ge 0\}$
(\texttt{performance.py}). For a joint ground-motion threshold condition with
per-output thresholds $\gamma_k$, define the per-event, per-output log
margin
\begin{equation}
r_{i,k} = \log\!\left(\frac{Y_{i,k}}{\gamma_k}\right),
\end{equation}
an event score aggregating across outputs $k=1,\dots,K$,
\begin{equation}
\text{event score}_i =
\begin{cases}
\min_k r_{i,k} & \text{mode = ``all'' (joint exceedance)}\\
\max_k r_{i,k} & \text{mode = ``any'' (union exceedance)}\\
r_{i,(m)} & \text{mode = ``$k$-of-$n$'', $m$-th largest margin,}
\end{cases}
\end{equation}
and the sequence-level performance as the maximum over events,
\begin{equation}
g(x) = \max_{i=1,\dots,N} \text{event score}_i .
\end{equation}
The failure event is then $F = \{x : g(x) \ge 0\}$, i.e. some event in the
sequence exceeds the threshold condition. General system-level conditions can
be supplied as an arbitrary callable of the full sample.

\subsection{Smoothed indicator and unnormalized target}

Direct use of the indicator $\mathbb 1_F(x)$ is not differentiable, so a
smooth penalty is used in place of a hard indicator during training
(\texttt{target.py: PenaltyFunction}),
\begin{equation}
\rho_\alpha(x) = \exp\big(-\alpha\,\mathrm{ReLU}(0 - g(x))\big)
\quad\Longleftrightarrow\quad
\log \rho_\alpha(x) = -\alpha\,\mathrm{ReLU}(-g(x)),
\label{eq:penalty}
\end{equation}
with steepness $\alpha > 0$. As $\alpha \to \infty$, $\rho_\alpha(x) \to
\mathbb 1_F(x)$; $\rho_\alpha(x) = 1$ whenever $g(x) \ge 0$ and decays
exponentially for $g(x) < 0$.

The unnormalized rare-event target density is
\begin{equation}
h(x) = p(x)\,\rho_\alpha(x),
\qquad
\log h(x) = \log p(x) + \log \rho_\alpha(x),
\label{eq:target}
\end{equation}
computed by \texttt{target.py: RareEventTargetDensity}. The (unknown)
normalizing constant of $h$ is closely related to the failure probability of
interest,
\begin{equation}
\Prob(F) = \E_{p}\big[\mathbb 1_F(X)\big] = \int \mathbb 1_F(x)\,p(x)\,dx
\approx \int \rho_\alpha(x)\,p(x)\,dx = \int h(x)\,dx
\end{equation}
in the limit $\alpha \to \infty$; during training $\alpha$ is annealed
upward (via a progressive threshold schedule, Section~\ref{sec:progressive})
so that $h$ increasingly concentrates on $F$.

\section{Autoregressive normalizing-flow proposal}

\subsection{Structure}

The trainable proposal $q_\theta(x)$ is an event-by-event autoregressive
normalizing flow (\texttt{flow.py: AutoregressiveEventFlow}) that mirrors the
generative structure of $p(x)$ but replaces every conditional Gaussian by a
learned Gaussian whose mean and scale are outputs of a neural network
conditioned on the full sequence history through a recurrent hidden state.

Within event $i$, variables are ordered
\begin{equation}
v_i = \big(\dt_i,\; \text{remaining source variables},\; \bm\epsilon_i\big)
\in \R^{m}, \qquad m = |\text{source vars}| + K,
\end{equation}
and are generated autoregressively \emph{within} the event and \emph{across}
events: variable $v_{i,j}$ is conditioned on the recurrent hidden state
$h_{i-1}$ (summarizing all previous events), the current/remaining exposure
time, and all previously generated raw variables $v_{i,1:j-1}$ of the same
event,
\begin{equation}
\big(\mu_{i,j}(\theta),\, \log\sigma_{i,j}(\theta)\big)
=
\mathrm{MLP}_j\big(h_{i-1},\, t_{i-1}/T,\, (T-t_{i-1})/T,\, v_{i,1:j-1}\big),
\label{eq:flow-params}
\end{equation}
implemented as one small conditional MLP per variable position
(\texttt{ConditionalMLP}), with $\log\sigma$ squashed into
$[\log\sigma_{\min}, \log\sigma_{\max}]$ via a sigmoid for numerical
stability. The raw (latent) variable is drawn
\begin{equation}
z_{i,j} \sim \Normal(0,1), \qquad
\tilde v_{i,j} = \mu_{i,j}(\theta) + \sigma_{i,j}(\theta)\, z_{i,j},
\end{equation}
and mapped through the same physical transform $T_j$ used by the source
model (identity for residual variables, since they are already standard
Gaussian in the physical model):
\begin{equation}
v_{i,j} = T_j(\tilde v_{i,j}).
\end{equation}

The event-level flow log-density (change of variables) is
\begin{equation}
\log q_\theta(v_i \mid h_{i-1})
=
\sum_{j=1}^m \Big[\Normal\big(\tilde v_{i,j}; \mu_{i,j}, \sigma_{i,j}^2\big)
- \log\lvert T_j'(\tilde v_{i,j})\rvert \Big].
\label{eq:flow-event-density}
\end{equation}

\subsection{Recurrent state update and stopping}

After event $i$ is generated, the hidden state is updated by a GRU cell
consuming the raw event vector and normalized event time,
\begin{equation}
h_i = \mathrm{GRUCell}\big([\tilde v_i,\; t_i/T],\; h_{i-1}\big),
\qquad h_0 = h_{\mathrm{init}} \text{ (learned)}.
\end{equation}
Sequence generation stops the first time a proposed event time
$t_{i-1} + \dt_i$ exceeds $T$; as in the physical model, the flow's
log-density accrues a matching stopping log-probability using the same
Gaussian time-increment head evaluated as a normal log-survival function
(Eq.~\eqref{eq:survival}, with $\mu,\sigma$ from Eq.~\eqref{eq:flow-params}
instead of the physical transition). The complete sequence log-density under
the proposal is
\begin{equation}
\log q_\theta(x)
=
\sum_{i=1}^N \log q_\theta(v_i \mid h_{i-1})
\;+\;
\log \Prob_\theta(\dt_{N+1} > T - t_N \mid h_N).
\label{eq:flow-density}
\end{equation}
This is a valid joint density over variable-length sequences: it is a
triangular (autoregressive) transform of an underlying standard-normal base
sample, so $q_\theta$ is normalized and samples/densities are both tractable
(\texttt{sample\_one} / \texttt{log\_prob\_one}).

\section{Learning problem}

\subsection{Objective: forward KL / weighted maximum likelihood}

The proposal is trained to approximate the (self-normalized) rare-event
target $h(x)/Z$, $Z = \int h$, by minimizing the forward Kullback--Leibler
divergence
\begin{equation}
\theta^\star
=
\arg\min_\theta \; D_{\mathrm{KL}}\!\left(\frac{h}{Z} \,\Big\|\, q_\theta\right)
=
\arg\max_\theta \; \E_{x \sim h/Z}\big[\log q_\theta(x)\big] + \text{const.}
\label{eq:kl}
\end{equation}
Since samples from $h/Z$ are not directly available, a self-normalized
importance-sampling / cross-entropy estimator is used
(\texttt{estimators.py: cross\_entropy\_update},
\texttt{training.py: defensive\_mixture\_update}). Draws are taken from a
\emph{defensive mixture} proposal that mixes the current flow with the
original physical model,
\begin{equation}
r_\theta(x) = (1-\varepsilon)\, q_\theta(x) + \varepsilon\, p(x),
\qquad \varepsilon \in (0,1),
\label{eq:defensive-mixture}
\end{equation}
which guarantees non-vanishing support under the physical model even while
$q_\theta$ is still adapting. For $M$ samples $x^{(1)},\dots,x^{(M)} \sim
r_\theta$, tempered self-normalized importance weights are formed,
\begin{equation}
w^{(m)}
=
\frac{\big(h(x^{(m)})/r_\theta(x^{(m)})\big)^{1/\tau}}
{\sum_{m'=1}^M \big(h(x^{(m')})/r_\theta(x^{(m')})\big)^{1/\tau}}
=
\mathrm{softmax}_m\!\left(\frac{\log h(x^{(m)}) - \log r_\theta(x^{(m)})}{\tau}\right),
\label{eq:weights}
\end{equation}
with temperature $\tau > 0$ smoothing the weight distribution, and the
weighted cross-entropy loss minimized by stochastic gradient (Adam) is
\begin{equation}
\mathcal L(\theta) = -\sum_{m=1}^M w^{(m)} \log q_\theta(x^{(m)}),
\label{eq:loss}
\end{equation}
which is an unbiased (in the large-$M$ limit) stochastic estimate of the
gradient of the KL objective in Eq.~\eqref{eq:kl}. Weights are computed under
\texttt{torch.no\_grad()} (self-normalized IS is treated as fixed) while only
$\log q_\theta(x^{(m)})$ carries gradients.

\subsection{Pretraining}

Before any rare-event adaptation, the flow is warm-started by ordinary
maximum likelihood on samples from the original physical model $p$
(\texttt{training.py: pretrain\_on\_original\_model}),
\begin{equation}
\theta_0 = \arg\min_\theta \; -\frac1M \sum_{m=1}^M \log q_\theta(x^{(m)}),
\qquad x^{(m)} \overset{\text{iid}}{\sim} p .
\end{equation}
This matches $q_\theta \approx p$ before any rare-event penalty is applied,
providing a stable starting point for cross-entropy adaptation.

\subsection{Progressive threshold schedule}
\label{sec:progressive}

Directly training against a very rare threshold $\gamma$ from initialization
is unstable (near-zero effective sample size). Instead, thresholds
$\gamma_1 < \gamma_2 < \dots < \gamma_L$ are visited in increasing order
(\texttt{training.py: ProgressiveThresholdTrainer.train\_schedule}), each
paired with a penalty steepness $\alpha_\ell$; at each stage the flow and
Adam optimizer state from the previous, easier threshold are retained as the
initialization,
\begin{equation}
\theta_\ell^{(0)} = \theta_{\ell-1}^{\star},
\qquad
\theta_\ell^{\star} = \arg\min_\theta \mathcal L_\ell(\theta),
\end{equation}
where $\mathcal L_\ell$ uses target $h_\ell(x) = p(x)\rho_{\alpha_\ell}(x)$
with threshold $\gamma_\ell$ substituted into the performance function. This
continuation strategy lets a proposal trained for an easier (more probable)
threshold serve as a good initialization for a rarer one, analogous to
sequential/adaptive importance sampling and cross-entropy methods for rare
events.

\subsection{Diagnostics}

Every update reports importance-weight diagnostics used to monitor learning
health: the maximum normalized weight $\max_m w^{(m)}$, the weight (Shannon)
entropy $-\sum_m w^{(m)}\log w^{(m)}$, and the weighted effective sample size
\begin{equation}
\mathrm{ESS} = \frac{1}{\sum_{m=1}^M (w^{(m)})^2},
\end{equation}
which should remain a non-trivial fraction of $M$ for a well-adapted
proposal.

\section{Sample generation summary}

Generating one complete simulated sample proceeds as follows.

\begin{enumerate}[label=\arabic*.]
\item \textbf{Initialize} time $t_0 = 0$ and the source state at a reference
event $\bs_0$ (physical model) or recurrent hidden state $h_0$ (flow).
\item \textbf{Loop over events} $i = 1, 2, \dots$:
  \begin{enumerate}[label=(\alph*)]
  \item Draw the latent Gaussian vector for the next event conditional on the
  previous state (Eq.~\eqref{eq:latent-transition} for $p$, or
  Eq.~\eqref{eq:flow-params} autoregressively coordinate-by-coordinate for
  $q_\theta$), including the time increment $\dt_i$ first.
  \item Map to physical units via the invertible transforms $T_j$.
  \item If $t_{i-1} + \dt_i > T$: stop, record the stopping (survival)
  log-probability, and terminate the sequence with $N = i-1$ events.
  \item Otherwise accept event $i$, set $t_i = t_{i-1}+\dt_i$, build ground-
  motion features $\bm f_i$ from the event and site, draw standard-normal
  residuals $\bm\epsilon_i$ (or their flow equivalents) and compute
  $\bm Y_i$ via Eq.~\eqref{eq:gmm}, and update the conditioning state/hidden
  state for the next event.
  \end{enumerate}
\item \textbf{Return} the full sample $x = (\mathcal X_N, \{\bm Y_i\})$
together with its log-density (Eq.~\eqref{eq:original-density} under $p$, or
Eq.~\eqref{eq:flow-density} under $q_\theta$).
\end{enumerate}

Because $q_\theta$ shares the same variable ordering, transforms, and
autoregressive/Markov factorization structure as $p$, it can generate
sequences of arbitrary (and different) length $N$ while remaining
absolutely continuous with respect to $p$ almost everywhere, which is what
makes it usable as an importance-sampling proposal.

\section{Rare-event probability estimation}

\subsection{Naive Monte Carlo}

Direct simulation from the physical model gives an unbiased but typically
very high-variance estimator for rare $F$ (\texttt{estimators.py:
estimate\_probability\_naive\_monte\_carlo}),
\begin{equation}
\hat P_{\mathrm{MC}} = \frac1M \sum_{m=1}^M \mathbb 1_F(x^{(m)}),
\qquad x^{(m)} \overset{\text{iid}}{\sim} p,
\end{equation}
with a Wilson binomial confidence interval, appropriate when the exceedance
count may be small or zero.

\subsection{Flow-only importance sampling}

Using the trained flow directly as an importance proposal
(\texttt{estimate\_probability\_flow\_only}),
\begin{equation}
\hat P_{q} = \frac1M \sum_{m=1}^M \mathbb 1_F(x^{(m)})\,
\frac{p(x^{(m)})}{q_\theta(x^{(m)})},
\qquad x^{(m)} \overset{\text{iid}}{\sim} q_\theta .
\end{equation}
This is unbiased provided $q_\theta$'s support covers $F$, but is used mainly
as a diagnostic since a poorly-adapted or mode-collapsed flow can miss
support entirely.

\subsection{Defensive-mixture importance sampling}

The safeguarded, final estimator draws from the same defensive mixture used
in training (Eq.~\eqref{eq:defensive-mixture}),
\begin{equation}
\hat P_{r} = \frac1M \sum_{m=1}^M \mathbb 1_F(x^{(m)})\,
\frac{p(x^{(m)})}{r_\theta(x^{(m)})},
\qquad x^{(m)} \sim r_\theta,
\end{equation}
which guarantees a bounded importance weight $p(x)/r_\theta(x) \le
p(x)/(\varepsilon\, p(x)) = 1/\varepsilon$ whenever the sample is drawn from
the physical component, preventing the unbounded-variance failure mode of a
pure learned-proposal estimator. All three estimators return standard error,
a confidence interval, and the importance-weight effective sample size as
diagnostics of estimator reliability.

\end{document}
